% Laget automatisk av elevark/lag_elevark.py fra sim/impuls.html. Endringer her blir overskrevet.
\documentclass[11pt]{article}
\usepackage{elevark}
\begin{document}
\elevarkhode{Fysikk 1 · Bevegelsesmengde og støt}{Myk landing}{https://fysikklab.no/sim/impuls.html}{Simuleringen}

Et egg på 60~g slippes fra en høyde og lander på betong, en treplate, gress, skumgummi eller en dyne. Hvorfor knuser egget på betong, men ikke på en dyne? Bytt så til en sprettball med samme masse: Hvorfor får ballen større endring i bevegelsesmengde enn egget når den spretter opp igjen?

\section*{Oppgaver}

Newtons 2. lov kan skrives som $ΣF \cdot Δt = Δp$. Produktet $F \cdot Δt$ kalles impuls, og det er arealet under grafen for $F$ mot $t$.

\begin{oppgaveliste}
\oppgave \fase{gjett}{Gjett}Egget slippes fra 1,0~m. Er endringen i bevegelsesmengde større når egget lander på betong enn når det lander på en dyne? Er kraften større?
\svar
\oppgave \fase{test}{Test}Slipp egget fra 1,0~m på alle de fem underlagene. Sammenlign bremsetiden, $Δp$ og den største kraften. Hva er likt, og hva er forskjellig?
\svar
\oppgave \fase{forklar}{Forklar}Grafene har svært ulik form. Hva kan du si om arealet under dem fram til egget står stille? Bruk $ΣF \cdot Δt = Δp$ til å forklare hvorfor egget knuser på det harde underlaget.
\svar
\oppgave \fase{test}{Regn}Regn ut farten rett før egget lander fra 1,0~m, og bevegelsesmengden det har da. Egget lander på skumgummi og bremses på 17~ms. Regn ut middelkraften fra underlaget. Husk at tyngden også virker.
\svar
\oppgave \fase{test}{Undersøk}Hvor høyt kan du slippe egget på skumgummi før det knuser? Hva skjer med den største kraften når du dobler høyden? Hvorfor blir den ikke dobbelt så stor?
\svar
\oppgave \fase{forklar}{Forklar}Bruk det samme resonnementet på kollisjonsputer og knasesoner i biler, hjelmer, sklimatter og at du bøyer knærne når du hopper ned fra en høyde.
\svar
\oppgave \fase{gjett}{Gjett}Velg sprettball og betong, og slipp fra 1,0~m. Ballen har samme masse og samme fart før landingen som egget. Blir $Δp$ større, mindre eller like stor som for egget, som blir liggende? Begrunn svaret før du slipper.
\svar
\oppgave \fase{test}{Regn}Les av farten ballen spretter opp med. Velg positiv retning oppover, og regn ut $p$ før, $p$ etter og $Δp = p$ etter − $p$ før. Sjekk med pilene og tallene i simuleringen. Hvorfor er $Δp$ nesten dobbelt så stor som for egget?
\svar
\oppgave \fase{forklar}{Forklar}Sammenlign arealet under grafen for egget og for ballen på betong. Hva viser den lyseste delen av arealet for ballen? Ballen har likevel mindre største kraft enn egget. Hvordan kan det ha seg?
\svar
\oppgave \fase{test}{Undersøk}Slipp ballen fra 1,0~m på alle underlagene. På hvilket underlag blir $Δp$ størst, og på hvilket blir den minst? Hvilke verdier kan $Δp$ ligge mellom?
\svar
\oppgave \fase{utfordring}{Utfordring}Eggslippkonkurranse: Lag en innpakning av papir, sugerør og teip som gjør at et egg tåler å slippes fra 3~m. Hvilken bremsetid må innpakningen gi hvis egget skal tåle 30~N?
\svar
\end{oppgaveliste}

\end{document}
